\documentclass[11pt]{article}
\usepackage{amsmath, amssymb, physics}
\usepackage[margin=1in]{geometry}
\usepackage{graphicx}
\usepackage{booktabs}
\usepackage{caption}
\usepackage{amsmath}
\setlength{\parindent}{0pt}
\setlength{\parskip}{6pt}

\newenvironment{solution}
{\par\noindent\textbf{Solution:}\par}
{\vspace{1em}}

\title{Homework Assignment No.1 \\ EE520/EE459 \\ Yuyao Huang}
\date{}
\begin{document}
\maketitle

\textbf{Total: 70 points} \\
\textbf{Submit: a single PDF file of all your work}

\section*{Problem 1 (10 points)}
From the definition, write down the Bloch sphere coordinates $(\theta,\phi)$ for the points of intercept on the $x,y,z$ axes by the Bloch sphere:
\[
(x,y,z)=(+1,0,0)=\ket{+},\;
(-1,0,0)=\ket{-},\;
(0,+1,0),\;
(0,-1,0),\;
(0,0,+1)=\ket{0},\;
(0,0,-1)=\ket{1}.
\]
Sketch the Bloch sphere and indicate the resulting states in terms of $\ket{0}, \ket{1}$ and their superposition.

\begin{solution}

The question gives the key information indicating that six special points on the Bloch sphere correspond to specific quantum states.

\begin{itemize}
\item $x$-axis $\rightarrow$ Hadamard phase basis ($\ket{+}$, $\ket{-}$),
\item $y$-axis $\rightarrow$ complex basis ($\ket{+i}$, $\ket{-i}$),
\item $z$-axis $\rightarrow$ computational basis ($\ket{0}$, $\ket{1}$).
\end{itemize}

So this information tells us there are three axes with positive and negative directions on the Bloch sphere.

We now determine the quantum state corresponding to each axis intercept by solving the eigenvalue equations of the Pauli operators. 
Let an arbitrary single–qubit state be written in the computational basis as
\[
\ket{\psi}=a\ket{0}+b\ket{1}=\begin{bmatrix}a\\ b\end{bmatrix}, 
\qquad a,b\in\mathbb{C},
\]
with normalization $|a|^2+|b|^2=1$.

The Pauli matrices are
\[
X=\begin{bmatrix}0&1\\1&0\end{bmatrix},\qquad
Y=\begin{bmatrix}0&-i\\ i&0\end{bmatrix},\qquad
Z=\begin{bmatrix}1&0\\0&-1\end{bmatrix}.
\]

\textbf{Z-axis intercepts.}

For the north pole $(0,0,+1)$ we solve $Z\ket{\psi}=+\ket{\psi}$:
\[
\begin{bmatrix}1&0\\0&-1\end{bmatrix}
\begin{bmatrix}a\\ b\end{bmatrix}
=
\begin{bmatrix}a\\ b\end{bmatrix}
\Rightarrow
\begin{bmatrix}a\\ -b\end{bmatrix}
=
\begin{bmatrix}a\\ b\end{bmatrix}
\Rightarrow b=0.
\]
After normalization, $\ket{\psi}=\ket{0}$.

For the south pole $(0,0,-1)$ we solve $Z\ket{\psi}=-\ket{\psi}$:
\[
\begin{bmatrix}a\\ -b\end{bmatrix}
=
\begin{bmatrix}-a\\ -b\end{bmatrix}
\Rightarrow a=0,
\]
which gives $\ket{\psi}=\ket{1}$.

\textbf{X-axis intercepts.}

For $(+1,0,0)$ we solve $X\ket{\psi}=+\ket{\psi}$:
\[
\begin{bmatrix}0&1\\1&0\end{bmatrix}
\begin{bmatrix}a\\ b\end{bmatrix}
=
\begin{bmatrix}a\\ b\end{bmatrix}
\Rightarrow
\begin{bmatrix}b\\ a\end{bmatrix}
=
\begin{bmatrix}a\\ b\end{bmatrix}
\Rightarrow b=a.
\]
Normalization gives
\[
\ket{+}=\frac{1}{\sqrt2}(\ket{0}+\ket{1}).
\]

For $(-1,0,0)$ we solve $X\ket{\psi}=-\ket{\psi}$:
\[
\begin{bmatrix}b\\ a\end{bmatrix}
=
\begin{bmatrix}-a\\ -b\end{bmatrix}
\Rightarrow b=-a,
\]
yielding
\[
\ket{-}=\frac{1}{\sqrt2}(\ket{0}-\ket{1}).
\]

\textbf{Y-axis intercepts.}

For $(0,+1,0)$ we solve $Y\ket{\psi}=+\ket{\psi}$:
\[
\begin{bmatrix}0&-i\\ i&0\end{bmatrix}
\begin{bmatrix}a\\ b\end{bmatrix}
=
\begin{bmatrix}a\\ b\end{bmatrix}
\Rightarrow
\begin{bmatrix}-ib\\ ia\end{bmatrix}
=
\begin{bmatrix}a\\ b\end{bmatrix}.
\]
This gives $b=ia$, and after normalization,
\[
\ket{+i}=\frac{1}{\sqrt2}(\ket{0}+i\ket{1}).
\]

For $(0,-1,0)$ we solve $Y\ket{\psi}=-\ket{\psi}$:
\[
\begin{bmatrix}-ib\\ ia\end{bmatrix}
=
\begin{bmatrix}-a\\ -b\end{bmatrix}
\Rightarrow b=-ia,
\]
which gives
\[
\ket{-i}=\frac{1}{\sqrt2}(\ket{0}-i\ket{1}).
\]

\textbf{Bloch sphere angles.}

Any pure qubit state can be written as
\[
\ket{\psi}=\cos\frac{q}{2}\ket{0}+e^{ij}\sin\frac{q}{2}\ket{1}.
\]
From the amplitudes of each state we obtain

\[
\boxed{\begin{aligned}
(0,0,+1)&:\ \ket{0} \Rightarrow (q,j)=(0,0),\\
(0,0,-1)&:\ \ket{1} \Rightarrow (q,j)=(\pi,0),\\
(+1,0,0)&:\ \ket{+} \Rightarrow (q,j)=\left(\frac{\pi}{2},0\right),\\
(-1,0,0)&:\ \ket{-} \Rightarrow (q,j)=\left(\frac{\pi}{2},\pi\right),\\
(0,+1,0)&:\ \ket{+i} \Rightarrow (q,j)=\left(\frac{\pi}{2},\frac{\pi}{2}\right),\\
(0,-1,0)&:\ \ket{-i} \Rightarrow (q,j)=\left(\frac{\pi}{2},\frac{3\pi}{2}\right).
\end{aligned}}
\]

\begin{figure}[h]
\centering
\includegraphics[width=0.6\textwidth]{bloch_sphere.png}
\caption{Bloch sphere showing the six basis states at axis intercepts.}
\end{figure}

\end{solution}

\section*{Problem 2 (5 points)}
Using the matrix representation of the Hadamard gate $H$ and the single-qubit states, find
\[
H\ket{0},\quad H\ket{1},\quad H\ket{+},\quad H\ket{-}
\]
where
\[
\ket{\pm}=\frac{1}{\sqrt{2}}(\ket{0}\pm\ket{1}).
\]

\begin{solution}

\textbf{1. preparation, write all objects as "matrix/vector"}
\textbf{1.1. Computational basis}
At the standard condition, we have
\[
\ket{0}=\begin{bmatrix}1\\0\end{bmatrix},\quad
\ket{1}=\begin{bmatrix}0\\1\end{bmatrix}.
\]

\textbf{1.2. Hadamard gate}
The matrix representation of the Hadamard gate is
\[
H = \frac{1}{\sqrt{2}} \begin{bmatrix}1&1\\1&-1\end{bmatrix}.
\]

\textbf{1.3. Formation of $|+\rangle$, $|-\rangle$}
\[
\ket{+} = \frac{1}{\sqrt{2}} (\ket{0} + \ket{1}) = \frac{1}{\sqrt{2}} \left(\begin{bmatrix}1\\0\end{bmatrix} + \begin{bmatrix}0\\1\end{bmatrix}\right) = \frac{1}{\sqrt{2}} \begin{bmatrix}1\\1\end{bmatrix} 
\]
\[
\ket{-} = \frac{1}{\sqrt{2}} (\ket{0} - \ket{1}) = \frac{1}{\sqrt{2}} \left(\begin{bmatrix}1\\0\end{bmatrix} - \begin{bmatrix}0\\1\end{bmatrix}\right) = \frac{1}{\sqrt{2}} \begin{bmatrix}1\\-1\end{bmatrix}
\]

\textbf{2. $H|0\rangle$ computation}
\[
H\ket{0} = \left(\frac{1}{\sqrt{2}} \begin{bmatrix}1&1\\1&-1\end{bmatrix}\right) \begin{bmatrix}1\\0\end{bmatrix} = \frac{1}{\sqrt{2}} \left(\begin{bmatrix}1&1\\1&-1\end{bmatrix} \begin{bmatrix}1\\0\end{bmatrix}\right) = \frac{1}{\sqrt{2}} \begin{bmatrix}1 \cdot 1 + 1 \cdot 0\\1 \cdot 1 + (-1) \cdot 0\end{bmatrix} = \frac{1}{\sqrt{2}} \begin{bmatrix}1\\1\end{bmatrix} = \ket{+}
\]

\textbf{3. $H|1\rangle$ computation}
\[
H\ket{1} = \left(\frac{1}{\sqrt{2}} \begin{bmatrix}1&1\\1&-1\end{bmatrix}\right) \begin{bmatrix}0\\1\end{bmatrix} = \frac{1}{\sqrt{2}} \left(\begin{bmatrix}1&1\\1&-1\end{bmatrix} \begin{bmatrix}0\\1\end{bmatrix}\right) = \frac{1}{\sqrt{2}} \begin{bmatrix}1 \cdot 0 + 1 \cdot 1\\1 \cdot 0 + (-1) \cdot 1\end{bmatrix} = \frac{1}{\sqrt{2}} \begin{bmatrix}1\\-1\end{bmatrix} = \ket{-}
\]

\textbf{4. $H|+\rangle$ computation}
\[
H\ket{+} = \left(\frac{1}{\sqrt{2}} \begin{bmatrix}1&1\\1&-1\end{bmatrix}\right) \left(\frac{1}{\sqrt{2}} \begin{bmatrix}1\\1\end{bmatrix}\right) = \frac{1}{2} \begin{bmatrix}1&1\\1&-1\end{bmatrix} \begin{bmatrix}1\\1\end{bmatrix} = \frac{1}{2} \begin{bmatrix}1\cdot 1 + 1 \cdot 1\\1 \cdot 1 + (-1) \cdot 1\end{bmatrix} = \frac{1}{2} \begin{bmatrix}2\\0\end{bmatrix} = \begin{bmatrix}1\\0\end{bmatrix} = \ket{0}
\]

\textbf{5. $H|-\rangle$ computation}
\[
H\ket{-} = \left(\frac{1}{\sqrt{2}} \begin{bmatrix}1&1\\1&-1\end{bmatrix}\right) \left(\frac{1}{\sqrt{2}} \begin{bmatrix}1\\-1\end{bmatrix}\right) = \frac{1}{2} \begin{bmatrix}1&1\\1&-1\end{bmatrix} \begin{bmatrix}1\\-1\end{bmatrix} = \frac{1}{2} \begin{bmatrix}1\cdot 1 + 1 \cdot (-1)\\1 \cdot 1 + (-1) \cdot (-1)\end{bmatrix} = \frac{1}{2} \begin{bmatrix}0\\2\end{bmatrix} = \begin{bmatrix}0\\1\end{bmatrix} = \ket{1}
\]

\textbf{Results summary}
\[\boxed{H\ket{0} = \ket{+}}\]
\[\boxed{H\ket{1} = \ket{-}}\]
\[\boxed{H\ket{+} = \ket{0}}\]
\[\boxed{H\ket{-} = \ket{1}}\]

\end{solution}

\section*{Problem 3 (10 points)}
Using matrix representations evaluate
\[
X^2,\quad Y^2,\quad Z^2,\quad H^2,\quad S^2,\quad T^2.
\]
\begin{solution}

\textbf{1. preparation, write all objects as "matrix/vector"}
\textbf{1.1. Pauli matrices and gates}
\[
X = \begin{bmatrix}0&1\\1&0\end{bmatrix}, \quad Y = \begin{bmatrix}0&-i\\i&0\end{bmatrix}, \quad Z = \begin{bmatrix}1&0\\0&-1\end{bmatrix}
\]

\textbf{1.2. Hadamard gate}
\[
H = \frac{1}{\sqrt{2}}\begin{bmatrix}1&1\\1&-1\end{bmatrix}
\]

\textbf{1.3. Phase gates}
\[
S = \begin{bmatrix}1&0\\0&i\end{bmatrix}, \quad T = \begin{bmatrix}1&0\\0&e^{i\pi/4}\end{bmatrix}
\]

\textbf{2. computations}
\[
X^2 = XX = \begin{bmatrix}0&1\\1&0\end{bmatrix} \begin{bmatrix}0&1\\1&0\end{bmatrix} = \begin{bmatrix}0 \cdot 0 + 1 \cdot 1&0 \cdot 1 + 1 \cdot 0\\1 \cdot 0 + 0 \cdot 1&1 \cdot 1 + 0 \cdot 0\end{bmatrix} = \begin{bmatrix}1&0\\0&1\end{bmatrix} = \boxed{I}
\]

\[
Y^2 = YY = \begin{bmatrix}0&-i\\i&0\end{bmatrix} \begin{bmatrix}0&-i\\i&0\end{bmatrix} = \begin{bmatrix}0 \cdot 0 + (-i) \cdot i&0 \cdot (-i) + (-i) \cdot 0\\i \cdot 0 + 0 \cdot i&i \cdot (-i) + 0 \cdot 0\end{bmatrix} = \begin{bmatrix}1&0\\0&1\end{bmatrix} = \boxed{I} \text{ with } i^2 = -1
\]

\[
Z^2 = ZZ = \begin{bmatrix}1&0\\0&-1\end{bmatrix} \begin{bmatrix}1&0\\0&-1\end{bmatrix} = \begin{bmatrix}1 \cdot 1 + 0 \cdot 0&1 \cdot 0 + 0 \cdot (-1)\\0 \cdot 1 + (-1) \cdot 0&0 \cdot 0 + (-1) \cdot (-1)\end{bmatrix} = \begin{bmatrix}1&0\\0&1\end{bmatrix} = \boxed{I}
\]

\[
\begin{aligned}
H^2 &= HH \\
&= \left(\frac{1}{\sqrt{2}}
\begin{bmatrix}
1 & 1\\
1 & -1
\end{bmatrix}\right)
\left(\frac{1}{\sqrt{2}}
\begin{bmatrix}
1 & 1\\
1 & -1
\end{bmatrix}\right) \\
&= \frac{1}{2}
\begin{bmatrix}
1 & 1\\
1 & -1
\end{bmatrix}
\begin{bmatrix}
1 & 1\\
1 & -1
\end{bmatrix} \\
&= \frac{1}{2}
\begin{bmatrix}
1\cdot1 + 1\cdot1 & 1\cdot1 + 1\cdot(-1)\\
1\cdot1 + (-1)\cdot1 & 1\cdot1 + (-1)\cdot(-1)
\end{bmatrix} \\
&= \frac{1}{2}
\begin{bmatrix}
2 & 0\\
0 & 2
\end{bmatrix} \\
&=
\begin{bmatrix}
1 & 0\\
0 & 1
\end{bmatrix}
= \boxed{I}
\end{aligned}
\]

\[
S^2 = SS = \begin{bmatrix}1&0\\0&i\end{bmatrix} \begin{bmatrix}1&0\\0&i\end{bmatrix} = \begin{bmatrix}1 \cdot 1 + 0 \cdot 0&1 \cdot 0 + 0 \cdot i\\0 \cdot 1 + i \cdot 0&0 \cdot 0 + i \cdot i\end{bmatrix} = \begin{bmatrix}1&0\\0&-1\end{bmatrix} = \boxed{Z}
\]

\[
T^2 = TT = \begin{bmatrix}1&0\\0&e^{i\pi/4}\end{bmatrix} \begin{bmatrix}1&0\\0&e^{i\pi/4}\end{bmatrix} = \begin{bmatrix}1 \cdot 1 + 0 \cdot 0&1 \cdot 0 + 0 \cdot e^{i\pi/4}\\0 \cdot 1 + e^{i\pi/4} \cdot 0&0 \cdot 0 + e^{i\pi/4} \cdot e^{i\pi/4}\end{bmatrix} = \begin{bmatrix}1&0\\0&e^{i\pi/2}\end{bmatrix} = \boxed{S} \text{ with } e^{i\pi/2} = i
\]

\end{solution}

\section*{Problem 4 (15 points)}
Apply the following definitions:

\[
U \text{ is a unitary operator if } U^\dagger U = I
\]
\[
U \text{ is a Hermitian if } U^\dagger = U
\]

For each of the single qubit operators: $H, X, Y, Z, S, T,$ and $R_x(-\pi/2)$, determine:

\begin{itemize}
\item Is it unitary?
\item Is it Hermitian?
\end{itemize}

Hint: The rotation operator for $-90^\circ$ about the $x$-axis:
\[
R_n(\alpha)=\exp\left(-i\,\mathbf{n}\cdot\boldsymbol{\sigma}\,\frac{\alpha}{2}\right)
=\cos\left(\frac{\alpha}{2}\right)-i\,\mathbf{n}\cdot\boldsymbol{\sigma}\sin\left(\frac{\alpha}{2}\right)
\]
where $\mathbf{n}$ is the axis of rotation and $\boldsymbol{\sigma}=iX + jY + kZ$.

\begin{solution}

\textbf{1. Start from the standard rotation operator}

\textbf{1.1. Standard rotation operator}
\[
R_n(\alpha) = \exp\left(-i \frac{\alpha}{2} \mathbf{n} \cdot \boldsymbol{\sigma} \right)
\]
It represents the unitary operator that rotates a qubit state by angle $\alpha$ about the axis $\mathbf{n}$ on the Bloch sphere, which is responsible for the generation of the rotation gates and a matrix by an angle.
\begin{itemize}
\item $\mathbf{n}$ is a unit vector that defines the axis of rotation in 3D space, expressed as $\mathbf{n} = (n_x, n_y, n_z)$, such as, x-axis $(1,0,0)$, y-axis $(0,1,0)$, z-axis $(0,0,1)$.
\item $\alpha$ is rotation angle, unit is radians, such as $-\pi/2$ for $R_x(-\pi/2)$.
\end{itemize}

\textbf{1.2. Pauli vector}
\[
\boldsymbol{\sigma} = \hat{i}X + \hat{j}Y + \hat{k}Z \Rightarrow \mathbf{n} \cdot \boldsymbol{\sigma} = n_x \sigma_x + n_y \sigma_y + n_z \sigma_z
\]

\textbf{1.3. Use the standard closed form (Euler/Pauli identity)}

A key identity for Pauli matrices is: 
\[
R_n(\alpha)=\exp\left(-i\,\mathbf{n}\cdot\boldsymbol{\sigma}\,\frac{\alpha}{2}\right)
=\cos\left(\frac{\alpha}{2}\right)-i\,\mathbf{n}\cdot\boldsymbol{\sigma}\sin\left(\frac{\alpha}{2}\right)
\]
This is exactly the “standard Quantum Mechanics” expansion quoted.

\textbf{2. Specialize to rotation about the x-axis}

Rotation about the x-axis means the axis unit vector is
\[
\mathbf{n} = \hat{i} = (1,0,0)
\]

Then,
\[
\mathbf{n} \cdot \boldsymbol{\sigma} = \hat{i} \cdot (\hat{i}X + \hat{j}Y + \hat{k}Z) = (1,0,0) \cdot (\hat{i}X + \hat{j}Y + \hat{k}Z) = 1 \cdot X + 0 \cdot Y + 0 \cdot Z = X
\]

\textbf{3. Plug in $\alpha = -\pi/2$}
\[
R_x(-\pi/2) = \cos(\pi/4)I - i \sin(-\pi/4)X = \frac{1}{\sqrt{2}}I - i\left(-\frac{1}{\sqrt{2}} \right)X = \frac{1}{\sqrt{2}}(I + iX)
\]

Here is,
\[
I = \begin{bmatrix}1&0\\0&1\end{bmatrix}, \quad X = \begin{bmatrix}0&1\\1&0\end{bmatrix}
\]

\[
\Rightarrow I + iX = \begin{bmatrix}1&0\\0&1\end{bmatrix} + i\begin{bmatrix}0&1\\1&0\end{bmatrix} = \begin{bmatrix}1& i\\ i& 1\end{bmatrix}
\]

Therefore,
\[
R_x(-\pi/2) = \frac{1}{\sqrt{2}} \begin{bmatrix}1& i\\ i& 1\end{bmatrix}
\]

\end{solution}

\section*{Problem 5 (10 points)}
Show that the four Bell states $\ket{\beta_{cd}}$ form an orthonormal basis for a 4-dimensional vector space. That is, show:
\[
\langle \beta_{cd}|\beta_{c'd'} \rangle = \delta_{cc'}\delta_{dd'}
\]
where
\[
\delta_{ij}=
\begin{cases}
1 & i=j\\
0 & i\neq j
\end{cases}
\]

Hints: (1) calculate the inner products for $cd=00,01,10,11$; and do ten separate calculations: for 00, inner product with 00, 01, 10, 11; for 01, inner product with 01, 10, 11; for 10, with 10, 11; and for 11, with 11. (Check these examples out, as well as the main HW problem, for your own learning.)

Examples of orthonormal basis:
\[
\{\ket{0},\ket{1}\}\quad\text{for 2D because }\langle0|0\rangle=1=\langle1|1\rangle\text{ and }\langle0|1\rangle=0.
\]
\[
\{\ket{+},\ket{-}\}\quad\text{for 2D because }\langle+|+\rangle=1=\langle-|-\rangle\text{ and }\langle+|-\rangle=0.
\]
\[
\{\ket{00},\ket{01},\ket{10},\ket{11}\}\quad\text{for 4D because }\langle00|00\rangle=\langle01|01\rangle=\langle10|10\rangle=\langle11|11\rangle=1,
\]
and it's not hard to verify all other inner-product relations are summarized as $\langle ij|km\rangle = \delta_{ik}\delta_{jm}$. Check these three examples out, as well as the main HW problem, for your own learning.

\begin{solution}

\textbf{1.Define the Bell states (standard convention)}


\textbf{Conditions}
\[
\ket{00} = \begin{bmatrix}1\\0\\0\\0\end{bmatrix},\quad
\ket{01} = \begin{bmatrix}0\\1\\0\\0\end{bmatrix},\quad
\ket{10} = \begin{bmatrix}0\\0\\1\\0\end{bmatrix},\quad
\ket{11} = \begin{bmatrix}0\\0\\0\\1\end{bmatrix}
\]

We use the standard four Bell states (\textbf{maximally entangled states of 2 qubits}) labeled by $cd \in \{00,01,10,11\}$,
\[
\ket{\beta_{00}} = \ket{\Phi^+} = \frac{1}{\sqrt{2}}(\ket{00} + \ket{11}) = \frac{1}{\sqrt{2}} \left (\begin{bmatrix}1\\0\\0\\0\end{bmatrix} + \begin{bmatrix}0\\0\\0\\1\end{bmatrix} \right) = \frac{1}{\sqrt{2}} \begin{bmatrix}1\\0\\0\\1\end{bmatrix}
\]
\[\ket{\beta_{01}} = \ket{\Phi^-} = \frac{1}{\sqrt{2}}(\ket{00} - \ket{11}) = \frac{1}{\sqrt{2}} \left (\begin{bmatrix}1\\0\\0\\0\end{bmatrix} - \begin{bmatrix}0\\0\\0\\1\end{bmatrix} \right) = \frac{1}{\sqrt{2}} \begin{bmatrix}1\\0\\0\\-1\end{bmatrix}
\]
\[\ket{\beta_{10}} = \ket{\Psi^+} = \frac{1}{\sqrt{2}}(\ket{01} + \ket{10}) = \frac{1}{\sqrt{2}} \left (\begin{bmatrix}0\\1\\0\\0\end{bmatrix} + \begin{bmatrix}0\\0\\1\\0\end{bmatrix} \right) = \frac{1}{\sqrt{2}} \begin{bmatrix}0\\1\\1\\0\end{bmatrix}
\]
\[\ket{\beta_{11}} = \ket{\Psi^-} = \frac{1}{\sqrt{2}}(\ket{01} - \ket{10}) = \frac{1}{\sqrt{2}} \left (\begin{bmatrix}0\\1\\0\\0\end{bmatrix} - \begin{bmatrix}0\\0\\1\\0\end{bmatrix} \right) = \frac{1}{\sqrt{2}} \begin{bmatrix}0\\1\\-1\\0\end{bmatrix}
\]

\begin{table}[h]
\centering
\captionsetup{justification=raggedright,singlelinecheck=false,labelsep=period,labelfont=bf}
\caption{Bell states summary}
\begin{tabular}{cccccc}
\toprule
Label ($cd$) & Name & State formula & Measurement relation & Relative phase & Key intuition \\
\midrule
00 & $\ket{\Phi^+}$ & $\frac{1}{\sqrt{2}}(\ket{00}+\ket{11})$ & same (00 or 11) & $+$ & correlated, in-phase \\
01 & $\ket{\Phi^-}$ & $\frac{1}{\sqrt{2}}(\ket{00}-\ket{11})$ & same (00 or 11) & $-$ & correlated, out-of-phase \\
10 & $\ket{\Psi^+}$ & $\frac{1}{\sqrt{2}}(\ket{01}+\ket{10})$ & opposite (01 or 10) & $+$ & anti-correlated, in-phase \\
11 & $\ket{\Psi^-}$ & $\frac{1}{\sqrt{2}}(\ket{01}-\ket{10})$ & opposite (01 or 10) & $-$ & anti-correlated, out-of-phase \\
\bottomrule
\end{tabular}
\end{table}

\textbf{2. Hint 2 is allowed to use inner product}

\textbf{2.1. Orthnormality}
\[
\braket{ij}{cd} = \delta_{ic} \delta_{jd}
\]

\textbf{2.2. Do the 10 separate inner products (exactly as Hint-1 requests)}

(1) Group A: for 00, with 00, 01, 10, 11

\[\begin{aligned}
\braket{\beta_{00}} &= \left(\frac{1}{\sqrt{2}}(\bra{00} + \bra{11})\right) \left(\frac{1}{\sqrt{2}}(\ket{00} + \ket{11})\right) \\
&= \frac{1}{2} \begin{bmatrix}1&0&0&1\end{bmatrix} \begin{bmatrix}1\\0\\0\\1\end{bmatrix} \\
&= \frac{1}{2} (1 \cdot 1 + 0 \cdot 0 + 0 \cdot 0 + 1 \cdot 1) \\
&= 1
\end{aligned}\]

\[\begin{aligned}
\braket{\beta_{00}}{\beta_{01}} &= \left(\frac{1}{\sqrt{2}}(\bra{00} + \bra{11})\right) \left(\frac{1}{\sqrt{2}}(\ket{00} - \ket{11})\right) \\
&= \frac{1}{2} \begin{bmatrix}1&0&0&1\end{bmatrix} \begin{bmatrix}1\\0\\0\\-1\end{bmatrix} \\
&= \frac{1}{2} (1 \cdot 1 + 0 \cdot 0 + 0 \cdot 0 + 1 \cdot (-1)) \\
&= 0
\end{aligned}\]

\[\begin{aligned}
\braket{\beta_{00}}{\beta_{10}} &= \left(\frac{1}{\sqrt{2}}(\bra{00} + \bra{11})\right) \left(\frac{1}{\sqrt{2}}(\ket{01} + \ket{10})\right) \\
&= \frac{1}{2} \begin{bmatrix}1&0&0&1\end{bmatrix} \begin{bmatrix}0\\1\\1\\0\end{bmatrix} \\
&= \frac{1}{2} (1 \cdot 0 + 0 \cdot 1 + 0 \cdot 1 + 1 \cdot 0) \\
&= 0
\end{aligned}\]

\[\begin{aligned}
\braket{\beta_{00}}{\beta_{11}} &= \left(\frac{1}{\sqrt{2}}(\bra{00} + \bra{11})\right) \left(\frac{1}{\sqrt{2}}(\ket{01} - \ket{10})\right) \\
&= \frac{1}{2} \begin{bmatrix}1&0&0&1\end{bmatrix} \begin{bmatrix}0\\1\\-1\\0\end{bmatrix} \\
&= \frac{1}{2} (1 \cdot 0 + 0 \cdot 1 + 0 \cdot (-1) + 1 \cdot 0) \\
&= 0
\end{aligned}\]

(2) Group B: for 01, with 01, 10, 11

\[\begin{aligned}
\braket{\beta_{01}}{\beta_{01}} &= \left(\frac{1}{\sqrt{2}}(\bra{00} - \bra{11})\right) \left(\frac{1}{\sqrt{2}}(\ket{00} - \ket{11})\right) \\
&= \frac{1}{2} \begin{bmatrix}1&0&0&-1\end{bmatrix} \begin{bmatrix}1\\0\\0\\-1\end{bmatrix} \\
&= \frac{1}{2} (1 \cdot 1 + 0 \cdot 0 + 0 \cdot 0 + (-1) \cdot (-1)) \\
&= 1
\end{aligned}\]

\[\begin{aligned}
\braket{\beta_{01}}{\beta_{10}} &= \left(\frac{1}{\sqrt{2}}(\bra{00} - \bra{11})\right) \left(\frac{1}{\sqrt{2}}(\ket{01} + \ket{10})\right) \\
&= \frac{1}{2} \begin{bmatrix}1&0&0&-1\end{bmatrix} \begin{bmatrix}0\\1\\1\\0\end{bmatrix} \\
&= \frac{1}{2} (1 \cdot 0 + 0 \cdot 1 + 0 \cdot 1 + (-1) \cdot 0) \\
&= 0
\end{aligned}\]

\[\begin{aligned}
\braket{\beta_{01}}{\beta_{11}} &= \left(\frac{1}{\sqrt{2}}(\bra{00} - \bra{11})\right) \left(\frac{1}{\sqrt{2}}(\ket{01} - \ket{10})\right) \\
&= \frac{1}{2} \begin{bmatrix}1&0&0&-1\end{bmatrix} \begin{bmatrix}0\\1\\-1\\0\end{bmatrix} \\
&= \frac{1}{2} (1 \cdot 0 + 0 \cdot 1 + 0 \cdot (-1) + (-1) \cdot 0) \\
&= 0
\end{aligned}\]

(3) Group C: for 10, with 10, 11

\[\begin{aligned}
\braket{\beta_{10}}{\beta_{10}} &= \left(\frac{1}{\sqrt{2}}(\bra{01} + \bra{10})\right) \left(\frac{1}{\sqrt{2}}(\ket{01} + \ket{10})\right) \\
&= \frac{1}{2} \begin{bmatrix}0&1&1&0\end{bmatrix} \begin{bmatrix}0\\1\\1\\0\end{bmatrix} \\
&= \frac{1}{2} (0 \cdot 0 + 1 \cdot 1 + 1 \cdot 1 + 0 \cdot 0) \\
&= 1
\end{aligned}\]

\[\begin{aligned}
\braket{\beta_{10}}{\beta_{11}} &= \left(\frac{1}{\sqrt{2}}(\bra{01} + \bra{10})\right) \left(\frac{1}{\sqrt{2}}(\ket{01} - \ket{10})\right) \\
&= \frac{1}{2} \begin{bmatrix}0&1&1&0\end{bmatrix} \begin{bmatrix}0\\1\\-1\\0\end{bmatrix} \\
&= \frac{1}{2} (0 \cdot 0 + 1 \cdot 1 + 1 \cdot (-1) + 0 \cdot 0) \\
&= 0
\end{aligned}\]

(4) Group D: for 11, with 11

\[\begin{aligned}
\braket{\beta_{11}}{\beta_{11}} &= \left(\frac{1}{\sqrt{2}}(\bra{01} - \bra{10})\right) \left(\frac{1}{\sqrt{2}}(\ket{01} - \ket{10})\right) \\
&= \frac{1}{2} \begin{bmatrix}0&1&-1&0\end{bmatrix} \begin{bmatrix}0\\1\\-1\\0\end{bmatrix} \\
&= \frac{1}{2} (0 \cdot 0 + 1 \cdot 1 + (-1) \cdot (-1) + 0 \cdot 0) \\
&= 1
\end{aligned}\]

\textbf{3. Results summary}

(1) Same Bell states:
\[\braket{\beta_{00}}{\beta_{00}} = 1, \quad \braket{\beta_{01}}{\beta_{01}} = 1, \quad \braket{\beta_{10}}{\beta_{10}} = 1, \quad \braket{\beta_{11}}{\beta_{11}} = 1\]

(2) Different Bell states:
\[\braket{\beta_{00}}{\beta_{01}} = 0, \quad \braket{\beta_{00}}{\beta_{10}} = 0, \quad \braket{\beta_{00}}{\beta_{11}} = 0, \quad \braket{\beta_{01}}{\beta_{10}} = 0, \quad \braket{\beta_{01}}{\beta_{11}} = 0, \quad \braket{\beta_{10}}{\beta_{11}} = 0\]

(3) Same labels imply same states, different labels imply different states.

\textbf{4. Conclusion}

From the results summary,
\begin{itemize}
    \item $\braket{\beta_{cd}}{\beta_{cd}} = 1$ for all $cd \in \{00, 01, 10, 11\}$ (normlized)
    \item $\braket{\beta_{cd}}{\beta_{c'd'}} = 0$ for all $cd \neq c'd'$ (orthogonal)
\end{itemize}

Therefore,
\[\boxed{\braket{\beta_{cd}}{\beta_{c'd'}} = \delta_{cc'}\delta_{dd'}}\]

\end{solution}

\section*{Problem 6 (10 points)}
For the CNOT operator, show it is self-inverse (\textit{i.e.}, $\text{CNOT}^2 = I$). From this self-inverse property, you can see that
\[
\text{CNOT}\ket{A}\ket{B}=\ket{A}\ket{A\oplus B}
\]
implies
\[
\text{CNOT}\ket{A}\ket{A\oplus B}=\ket{A}\ket{B}.
\]

\begin{solution}

\textbf{1. Definition of CNOT}
\[\text{CNOT}\ket{A}\ket{B} = \ket{A}\ket{A \oplus B} \text{ with } A,B \in \{0, 1\}\]
where $\oplus$ is the XOR operation.

\textbf{2. Show $\text{CNOT}^2 = I$}

Now we know,
\[
\text{CNOT}\ket{A}\ket{B} = \ket{A}\ket{A \oplus B}
\]

Applying CNOT again,
\[\text{CNOT}^2 \ket{A}\ket{B} = \text{CNOT} (\ket{A}\ket{A \oplus B}) = \ket{A}\ket{A \oplus (A \oplus B)}
\]

Now using the XOR identity,
\[\begin{aligned}
A \oplus A &= 0, \quad 0 \oplus B = B \\
\Rightarrow A \oplus (A \oplus B) &= (A \oplus A) \oplus B = 0 \oplus B = B
\end{aligned}\]

\[
\Rightarrow \ket{A}\ket{A \oplus (A \oplus B)} = \ket{A}\ket{B}
\]

Therefore,
\[\begin{aligned}
\text{CNOT}^2 \ket{A}\ket{B} &= \ket{A}\ket{B} \\
\Rightarrow \boxed{\text{CNOT}^2 = I}
\end{aligned}\]

\textbf{3. Deduce the “undo” statement}

Given
\[
\text{CNOT}\ket{A}\ket{B} = \ket{A}\ket{A \oplus B}
\]

Applying CNOT to both sides,
\[
\text{CNOT}(\text{CNOT}\ket{A}\ket{B}) = \text{CNOT}(\ket{A}\ket{A \oplus B})
\]

Left-hand side,
\[
\text{CNOT}^2 \ket{A}\ket{B} = \ket{A}\ket{B}
\]

Right-hand side,
\[
\text{CNOT}\ket{A}\ket{A \oplus B} = \ket{A}\ket{A \oplus (A \oplus B)} = \ket{A}\ket{B}
\]

Result summary,
\[
\boxed{\text{CNOT}\ket{A}\ket{A \oplus B} = \ket{A}\ket{B}}
\]

\end{solution}

\section*{Problem 7 (10 points)}
It is common to entangle separable states by applying $H$ followed by CNOT, and disentangle by applying CNOT followed by $H$.

Show that

\begin{figure}[h]
\centering
\includegraphics[width=0.9\textwidth]{problem7.png}
\end{figure}

\begin{solution}

This solution will use the above conditions and results from previous problems, such as the Hadamard gate transformations and the CNOT operation.

\textbf{1. Start from the definition will be helpful}

Hadamard,
\[
H\ket{x} = \frac{\ket{0} + (-1)^x \ket{1}}{\sqrt{2}}
\]

CNOT,
\[
\text{CNOT}\ket{A}\ket{B} = \ket{A}\ket{A \oplus B}
\]

Bell states,
\[
\ket{\beta_{00}} = \frac{\ket{00}+\ket{11}}{\sqrt{2}}, \quad \ket{\beta_{01}} = \frac{\ket{00}-\ket{11}}{\sqrt{2}}, \quad \ket{\beta_{10}} = \frac{\ket{01}+\ket{10}}{\sqrt{2}}, \quad \ket{\beta_{11}} = \frac{\ket{01}-\ket{10}}{\sqrt{2}}
\]

\textbf{2. Entangle: apply $H$ then CNOT to $\ket{x}\ket{y}$}

\textbf{2.1. Apply $H$ to the first qubit}
\[
(H \otimes I)\ket{x}\ket{y} = (H\ket{x}) \otimes \ket{y} = \frac{1}{\sqrt{2}}(\ket{0} + (-1)^x\ket{1})\ket{y} = \frac{1}{\sqrt{2}}(\ket{0}\ket{y} + (-1)^x\ket{1}\ket{y})
\]

\textbf{2.2. Apply CNOT to the resulting state}

Use CNOT definition,
\begin{itemize}
    \item $\text{CNOT}\ket{0}\ket{y} = \ket{0}\ket{0 \oplus y} = \ket{0}\ket{y}$
    \item $\text{CNOT}\ket{1}\ket{y} = \ket{1}\ket{1 \oplus y} = \ket{1}\ket{y \oplus 1}$
\end{itemize}
\[
\text{CNOT} \left( \frac{1}{\sqrt{2}}(\ket{0}\ket{y} + (-1)^x\ket{1}\ket{y}) \right) = \frac{1}{\sqrt{2}}(\ket{0}\ket{y} + (-1)^x\ket{1}\ket{1 \oplus y})
\]
That is the entangled state produced from $\ket{x}\ket{y}$.

\textbf{3. Match it to the four Bell states (do the four cases)}

Now check $y = 0$ and $y = 1$,

\textbf{case 1:} $y = 0$, then $y \oplus 1 = 1$, so
\[
\frac{1}{\sqrt{2}}(\ket{0}\ket{0} + (-1)^x\ket{1}\ket{1}) = \frac{1}{\sqrt{2}}(\ket{00} + (-1)^x\ket{11})
\]
\begin{itemize}
    \item If $x = 0$: $\frac{\ket{00} + \ket{11}}{\sqrt{2}} = \ket{\beta_{00}}$
    \item If $x = 1$: $\frac{\ket{00} - \ket{11}}{\sqrt{2}} = \ket{\beta_{01}}$
\end{itemize}
So for $y = 0$: output is $\ket{\beta_{00}}$ when $x = 0$, and $\ket{\beta_{01}}$ when $x = 1$.

\textbf{case 2:} $y = 1$, then $y \oplus 1 = 0$, so
\[
\frac{1}{\sqrt{2}}(\ket{0}\ket{1} + (-1)^x\ket{1}\ket{0}) = \frac{1}{\sqrt{2}}(\ket{01} + (-1)^x\ket{10})
\]
\begin{itemize}
    \item If $x = 0$: $\frac{\ket{01} + \ket{10}}{\sqrt{2}} = \ket{\beta_{10}}$
    \item If $x = 1$: $\frac{\ket{01} - \ket{10}}{\sqrt{2}} = \ket{\beta_{11}}$
\end{itemize}

Therefore, for all four $(x, y)$, the result is exactly the Bell state $\ket{\beta_{xy}}$:
\[
\textbf{CNOT}(H \otimes I)\ket{x}\ket{y} = \ket{\beta_{xy}}
\]

\textbf{4. Disentangle: apply CNOT then $H$ to the Bell state}

Now start from the Bell state we just got,
\[
\ket{\beta_{xy}} = \frac{1}{\sqrt{2}}(\ket{0}\ket{y} + (-1)^x \ket{1}\ket{1 \oplus y})
\]

\textbf{4.1. Apply CNOT}
\[
\ket{\beta_{xy}} = \frac{1}{\sqrt{2}}(\text{CNOT}\ket{0}\ket{y} + (-1)^x \text{CNOT}\ket{1}\ket{1 \oplus y})
\]

Compute each term,
\begin{itemize}
    \item $\text{CNOT}\ket{0}\ket{y} = \ket{0}\ket{y}$
    \item $\text{CNOT}\ket{1}\ket{1 \oplus y} = \ket{1}\ket{(1 \oplus y) \oplus 1} = \ket{1}\ket{y}$
\end{itemize}
Therefore,
\[
\text{CNOT}\ket{\beta_{xy}} = \frac{1}{\sqrt{2}}(\ket{0}\ket{y} + (-1)^x\ket{1}\ket{y}) = \left(\frac{\ket{0} + (-1)^x\ket{1}}{\sqrt{2}}\right)\ket{y}
\]

But $\frac{\ket{0} + (-1)^x\ket{1}}{\sqrt{2}} = H\ket{x}$. Hence,
\[
\text{CNOT}\ket{\beta_{xy}} = (H\ket{x})\ket{y}
\]

\textbf{4.2. Apply $H$ to the first qubit}
\[
(H \otimes I) \text{CNOT}\ket{\beta_{xy}} = (H(H\ket{x}))\ket{y} = (H^2\ket{x})\ket{y} = \ket{x}\ket{y}
\]
because $H^2 = I$, so,
\[\boxed{(H \otimes I) \text{CNOT}\ket{\beta_{xy}} = \ket{x}\ket{y}}
\]
That is precisely “disentangle by CNOT followed by $H$.”

\end{solution}

\end{document}

